\documentclass[aspectratio=169,11pt]{beamer}

\usetheme{Madrid}
\usecolortheme{default}
\usefonttheme{professionalfonts}
\setbeamertemplate{navigation symbols}{}
\setbeamertemplate{footline}[frame number]

\usepackage{amsmath, amssymb, booktabs, graphicx}

\title{Why Gender Matters For Labor Economics}
\subtitle{Gender Study --- Week 1}
\author{Teaching deck draft}
\date{}

\begin{document}

\begin{frame}
  \titlepage
\end{frame}

\begin{frame}{Central question}
\begin{itemize}
    \item What makes gender a core labor-economics object rather than a single wage-gap statistic?
    \item The answer starts from labor objects: employment, hours, earnings, wages, occupations, firms, authority, safety, care, and welfare.
    \item Week 1 discipline: name the outcome, state the comparison, identify the mechanism, and choose the right measurement window.
\end{itemize}
\end{frame}

\begin{frame}{Why this is a labor field}
\begin{itemize}
    \item Gender changes how workers allocate time, enter jobs, move across firms, receive pay, gain authority, face risk, and experience welfare.
    \item It connects Labor I objects: labor supply, family allocation, human capital, amenities, discrimination, and policy incidence.
    \item It connects Labor II objects: firms, search, wage-setting, contracts, regulation, internal labor markets, and worker voice.
    \item It is not only about wages, not only about women, and not only about discrimination.
\end{itemize}
\end{frame}

\begin{frame}{Core labor-market objects}
\small
\begin{tabular}{p{0.25\linewidth} p{0.60\linewidth}}
\toprule
Object & Labor question \\
\midrule
Employment & Who is attached to market work? \\
Hours and schedules & How much work is feasible, and when? \\
Wages and earnings & What is paid per hour, and what is earned over time? \\
Occupations and tasks & Where are workers allocated? \\
Firms and workplaces & Which firms hire, pay, promote, and retain? \\
Authority and promotion & Who controls resources, teams, and careers? \\
Safety and care & What risks and unpaid work shape the feasible set? \\
Welfare & What does work mean for utility, autonomy, and security? \\
\bottomrule
\end{tabular}
\end{frame}

\begin{frame}{Description starts with a gap}
\[
G_Y = \mathbb{E}[Y \mid F=1] - \mathbb{E}[Y \mid F=0]
\]
\begin{itemize}
    \item $Y$ can be employment, hours, earnings, wages, promotion, authority, safety, or welfare.
    \item The statistic is useful for magnitude, trends, and comparisons.
    \item It is not a mechanism: the same gap can reflect supply, care, sorting, firms, bargaining, law, safety, or discrimination.
\end{itemize}
\end{frame}

\begin{frame}{Why average gaps are incomplete}
\[
G = G^{prices} + G^{quantities} + G^{allocation}
\]
\small
\begin{tabular}{p{0.22\linewidth} p{0.66\linewidth}}
\toprule
Component & Interpretation \\
\midrule
Prices & pay-setting, wage premia, returns, bargaining, discrimination \\
Quantities & hours, employment, experience, interruptions, attachment \\
Allocation & occupations, tasks, firms, supervisors, schedules, career ladders \\
\bottomrule
\end{tabular}
\vspace{0.4cm}
\begin{itemize}
    \item Conceptual point: an earnings gap is a bundle until the mechanism is named.
\end{itemize}
\end{frame}

\begin{frame}{Mechanism classes}
\small
\begin{tabular}{p{0.25\linewidth} p{0.61\linewidth}}
\toprule
Class & Examples \\
\midrule
Worker-side & skills, expectations, search, commuting, amenities, risk \\
Household-side & care work, fertility, bargaining, family formation \\
Firm-side & hiring, evaluation, pay-setting, promotion, authority, networks \\
Institutional & law, enforcement, norms, legal rights, workplace protections \\
Policy & childcare, leave, taxation, quotas, transparency, safety regulation \\
\bottomrule
\end{tabular}
\vspace{0.4cm}
\begin{itemize}
    \item Later weeks unpack these channels one at a time.
\end{itemize}
\end{frame}

\begin{frame}{Measurement choices discipline claims}
\begin{itemize}
    \item Timing: entry, early career, family formation, mid-career, retirement.
    \item Family status: parents, nonparents, partnered workers, care responsibilities.
    \item Occupation and task: broad codes can hide authority, exposure, and schedules.
    \item Firm: workplace pay premia, promotion systems, supervisors, and climate.
    \item Legal environment: rights, restrictions, enforcement, eligibility, protections.
    \item Selection: wage data observe only workers with recorded jobs and pay.
\end{itemize}
\end{frame}

\begin{frame}{Lifecycle timing and child penalties}
\[
Y_{it} =
\sum_{k \neq -1}\beta_k \mathbf{1}\{t - T_i = k\}
+ \alpha_i + \gamma_t + \varepsilon_{it}
\]
\begin{itemize}
    \item $T_i$ indexes a family event such as first birth.
    \item Event time asks when gaps open, not only whether a cross-sectional gap exists.
    \item The outcome can be employment, hours, earnings, wages, occupation, or firm assignment.
    \item Week 2 turns this into care, fertility, and household allocation.
\end{itemize}
\end{frame}

\begin{frame}{Firms, careers, and authority}
\[
Y_{if} = \alpha_i + \psi_f + \delta F_i + \varepsilon_{if}
\]
\begin{itemize}
    \item Firms allocate jobs, schedules, clients, training, bonuses, promotions, and authority.
    \item Worker sorting and firm behavior can both shape observed gender gaps.
    \item Authority gaps may persist even when wage gaps look smaller.
    \item Later weeks move into hiring, promotion, pay-setting, transparency, and leadership pipelines.
\end{itemize}
\end{frame}

\begin{frame}{Safety, care, and welfare}
\begin{itemize}
    \item Safety changes the feasible set of jobs, commutes, hours, occupations, and workplaces.
    \item Care work changes market labor supply, career timing, schedule constraints, and policy incidence.
    \item Welfare is broader than earnings: flexibility, security, autonomy, dignity, exposure, and nonmarket time matter.
    \item Standard payroll data often miss these objects, so measurement has to expand.
\end{itemize}
\end{frame}

\begin{frame}{Legal and institutional dimensions}
\begin{itemize}
    \item Gendered laws affect access to work, mobility, assets, family decisions, safety, and workplace protections.
    \item Cross-country evidence is useful when it links legal objects to labor channels.
    \item Institutional measures must be interpreted carefully: the index is not itself the mechanism.
    \item Comparative settings help ask which mechanisms travel and which depend on development or labor-market structure.
\end{itemize}
\end{frame}

\begin{frame}{Why this is a high-return research field}
\begin{itemize}
    \item Remaining disparities are large in employment, hours, earnings, careers, authority, safety, care, and welfare.
    \item The field links labor supply, firms, contracts, occupations, institutions, policy, and inequality.
    \item Data are rich: administrative panels, linked worker-firm records, legal datasets, time use, surveys, audits, job postings.
    \item Methods travel: decompositions, event studies, linked data, field experiments, and comparative institutional designs.
    \item Frontier questions remain abundant enough for dissertation-scale projects.
\end{itemize}
\end{frame}

\begin{frame}{Week 1 lab}
\begin{itemize}
    \item Reproduce: synthetic country-year factbook inspired by gendered-laws evidence.
    \item Diagnose: outcome, comparison group, mechanism, level, and design type.
    \item Transfer: apply the same diagnostic logic to lifecycle child-penalty event time.
    \item Optional frontier prompt: how the question changes in worker-firm career data.
\end{itemize}
\end{frame}

\begin{frame}{Bridge to Week 2}
\begin{itemize}
    \item Week 1 gives the full map of labor objects and measurement choices.
    \item Week 2 zooms in on labor supply, care work, fertility, and household allocation.
    \item Central next question: how do care responsibilities and family formation alter employment, hours, earnings, job choice, and welfare over the lifecycle?
\end{itemize}
\end{frame}

\begin{frame}{Takeaways}
\begin{itemize}
    \item Gender is a core labor-economics field because it organizes work, pay, authority, risk, care, and welfare.
    \item Average gaps are necessary descriptions, but mechanisms require outcome choice, timing, and level of analysis.
    \item The course will move from objects to mechanisms to policy and frontier research design.
\end{itemize}
\end{frame}

\end{document}
